\begin{table}[t]
\tbl{\label{tab:basic:number:theory}}
{\begin{tcolorbox}[tabularx={p{3cm}|p{2cm}|p{6.3cm}|},enhanced,fonttitle=\bfseries\large,fontupper=\normalsize\sffamily,
colback=yellow!10!white,colframe=red!50!black,colbacktitle=blue!30!white,
coltitle=black,title=Basic Number Theory Commands]
\hline
Command & Arguments & Purpose \\
\hline
\hline
\verb'-primitive_root'  &  $p$  &  Computes a primitive root modulo $p$. \\
\hline
\verb'-smallest_primitive_root'  &  $p$  &  Computes the smallest primitive root modulo $p$. \\
\hline
\verb'-smallest_primitive_root_interval'  &  m1 m2  &  Computes the smallest primitive root modulo m for all $m$ with $m1 \le m \le m2$. \\
\hline
\verb'-number_of_primitive_roots_interval'  &  m1 m2  &   Counts the number of primitive roots modulo m for all $m$ with $m1 \le m \le m2$. \\
\hline
\verb'-inverse_mod'  &  $a$ $p$  &  Computes the modular inverse of $a$ modulo $p$, i.e. an integer $b$ with $ab \equiv 1$ mod $p$. \\
\hline
\verb'-extended_gcd'  &  $a$ $b$  &  Computes integers $g,$ $u,$ and $v$ such that $g = \gcd(a,b) = ua + vb.$ \\
\hline
\verb'-power_mod'  &  $a$ $k$ $n$  &  Compute the power $a^k$ mod $n$. \\
\hline
\verb'-discrete_log'  &  $y$ $a$ $m$  &  Computes $x$ such that $a^x \equiv y$ mod $m$. Employs a brute force algorithm. \\
\hline
\verb'-square_root'  &  $n$  &  computes $\lfloor \sqrt{a} \rfloor$ over the integers. \\
\hline
\verb'-square_root_mod'  &  $a$ $p$  &  Computes an integer $b$ such that $b^2 \equiv a$ mod $p$. Here, $p$ must be a prime number. \\
\hline
\verb'-all_square_roots_mod_n'  &  $a$ $n$  &  Computes all integers $b$ such that $b^2 \equiv a$ mod $n$. \\
\hline
\verb'-count_subprimitive'  &  Q\_max H\_max  &  Count subprimitive polynomials for $q \le Q\_max$ \\
\hline
\verb'-order_of_q_mod_n'  &  $q$ $n_{\min}$ $n_{\max}$  &  Computes the order ${\rm ord}(q,n)$ of $q$ modulo $n$ for all $n$ with $n_{\min} \le n \le n_{\max}$ for which $\gcd(n,q) = 1$. Also computes $\varphi(n)$ and $\varphi(n)/{\rm ord}(q,n)$. \\
\hline
\verb'-eulerfunction_interval'  &  n1 n2  &  Compute Euler's totient function $\varphi(n)$ for all n with $n1 \le n \le n2.$ \\
\hline
\verb'-jacobi'  &  $a$ $p$  &  Computes the Jacobi symbol $\left( \frac{a}{p} \right)$\\
\hline
\verb'-Chinese_remainders'  &  R M  &  Solves a system of congruences with remainders R and moduli M. R and M must be vectors whose labels are given. \\
\hline
\verb'-draw_mod_n'  &  option-string options &  Draws the integers modulo $n$ on a circle. For possible options, see Table~\ref{tab:draw:mod:n}. \\
\hline
\end{tcolorbox}}
\end{table}
